% keyence2019 B - KEYENCE String の手書きメモの再現と振り返り
% ビルド: tools/build_thinking.sh keyence2019/b --preview
\documentclass[a4paper]{article}
\usepackage[margin=1.2cm]{geometry}
\usepackage{handmemo}
\usepackage{hyperref}
\pagestyle{empty}

% 取り除く区間の「○」（原本では文字より少し大きい丸）
\newcommand{\hwo}{{\huge ○}}
% 原本の英字は方眼 2 マスほどの高さなので、このページの英字は一回り大きくする
\newcommand{\hwk}[1]{{\LARGE #1}}

\begin{document}

% ---------------------------------------------------------------- 1 枚目
\begin{memopage}{20260925\_174816.jpg}
  \hw(1.4,1.9){N日間旅行}
  \hw(5.7,1.9){長さNの文字列S}
  \redpen(10.6,1.3){\small 見出しの語は KEYENCE String の\\\small 問題文には無い}

  % 1 行目: ○key ? ence○
  \hw(1.0,4.75){\hwo} \hw(1.8,4.75){\hwk{key}}
  \hwunread(3.5,4.75){○かe}
  \hw(6.0,4.75){\hwk{ence}} \hw(7.8,4.75){\hwo}

  % 2 行目: ○k○ eyence○（eyence に下線）
  \hw(0.9,7.0){\hwo} \hw(1.8,7.0){\hwk{k}} \hw(2.6,7.0){\hwo}
  \hw(4.7,7.0){\hwk{eyence}} \hw(7.1,7.0){\hwo}
  \hwline(4.75,7.6)(7.0,7.5)

  % 3 行目: ○keyen○ce○（最後の○は少し上）
  \hw(0.7,9.0){\hwo} \hw(2.0,9.0){\hwk{keyen}} \hw(4.3,9.0){\hwo} \hw(5.3,9.0){\hwk{ce}}
  \hw(6.4,8.5){\hwo}

  % 4 行目: ○key enc○e○
  \hw(0.5,10.95){\hwo} \hw(2.0,10.95){\hwk{key}} \hw(3.6,10.95){\hwk{enc}} \hw(5.0,10.95){\hwo} \hw(6.1,10.95){\hwk{e}} \hw(6.9,10.95){\hwo}

  % 5 行目: ○ke○ yence○（yence に下線）
  \hw(0.9,13.2){\hwo} \hw(1.8,13.2){\hwk{ke}} \hw(2.7,13.2){\hwo}
  \hw(3.9,13.2){\hwk{yence}} \hw(6.1,13.2){\hwo}
  \hwline(3.9,13.8)(5.9,13.75)

  \redpen(9.0,5.6){取り除けるのは 1 か所だけ。\\S $=$ keyence[:k] $+$（何か）$+$ keyence[k:]\\分け目 k は 0〜7 の 8 通りしかない}
  \redpen(9.0,9.3){ここに書いた分け目は k=1, 5, 6, 2\\（1 行目は真ん中の字しだいで k=3）。\\k=0（○keyence）, 4, 7（keyence○）が無い}
  \redarrow(8.9,9.6)(7.8,10.7)

  \redpen(0.8,16.2){提出（1 回目で AC）は分け目 8 通りではなく、取り除く区間 [i, j) を\\全部試す O(N$^2$) だった。ただし {\hwlfont for j in range(i, l)} で j が l に届かず、\\末尾を取り除く形（k=7, keyence○）を試していない}
  \redpen(0.8,17.5){→ {\hwlfont keyencex} で NO を出す（正しくは YES）。テストにこの形が無く AC した}
  \redbox(0.6,15.2)(17.4,18.2)
\end{memopage}

% ---------------------------------------------------------------- 振り返り
\clearpage
\section*{keyence2019 B --- KEYENCE String}

\begin{itemize}
  \item 問題: \url{https://atcoder.jp/contests/keyence2019/tasks/keyence2019_b}
  \item 提出（2026-08-04, Python）:
    \href{https://atcoder.jp/contests/keyence2019/submissions/78117551}{AC}（1 回目で AC）
  \item 手書き原本: \texttt{Box/Photo Backup/手書き/atcoder/keyence2019/b/}（\texttt{20260925\_174816.jpg} の 1 枚。撮影 2026-09-25。前の 1 ページが再現。赤は後から入れた振り返り）
\end{itemize}

\subsection*{つまずき}
\begin{itemize}
  \item \textbf{メモ}: 取り除く区間（○）が keyence のどこに入るかを並べ、「前半 $+$ 後半に分ける位置は 8 通りしかない」形まで見えていた。
    ただし書いた分け目は 5 つで、$k=0$（先頭を取る）・$k=4$・$k=7$（末尾を取る）が無い。
  \item \textbf{提出}: メモの「分け方 8 通り」ではなく、取り除く区間 $[i,j)$ を $O(N^2)$ 通り全部試した（$N \leq 100$ なので間に合う）。
    ループが \verb|for j in range(i, l)| で $j = l$ を試さないため、末尾を取り除く形が抜けている。
    \verb|keyencex| や \verb|keyenceabc| で \verb|NO| を出す（正しくは \verb|YES|）。実際に動かして確かめた。
    テストケースにこの形が無く AC になった。\verb|range(i, l + 1)| にすれば通る。
\end{itemize}

\subsection*{正しい考え方}
$k = 0, \dots, 7$ について、$S$ が \verb|keyence[:k]| で始まり \verb|keyence[k:]| で終わるかを見る。$O(8N)$。
下の図は $S=$ \texttt{keyofscience}, $k=3$ の例。

\begin{center}
\begin{tikzpicture}[x=0.62cm,y=0.62cm]
  % 例: S = keyofscience, k = 3
  \foreach \c [count=\i] in {k,e,y,o,f,s,c,i,e,n,c,e} {\node[draw,minimum size=0.6cm,font=\ttfamily] at (\i,0) {\c};}
  \draw[decorate,decoration={brace,amplitude=5pt}] (0.6,0.55) -- (3.4,0.55) node[midway,above=6pt] {\small {\ttfamily keyence[:3]} で始まる};
  \draw[decorate,decoration={brace,amplitude=5pt}] (8.6,0.55) -- (12.4,0.55) node[midway,above=6pt] {\small {\ttfamily keyence[3:]} で終わる};
  \draw[redpen,line width=1.2pt] (3.6,-0.55) -- (8.4,-0.55);
  \node[redpen,below] at (6,-0.55) {\small ここ（{\ttfamily ofsci}）を 1 か所取り除く};
  \node[anchor=west] at (13.2,0) {\small $k=0$ なら先頭、$k=7$ なら末尾を取る};
\end{tikzpicture}
\end{center}

\begin{verbatim}
s = input()
t = "keyence"
print("YES" if any(s.startswith(t[:k]) and s.endswith(t[k:]) for k in range(8)) else "NO")
\end{verbatim}

\subsection*{次に活かすこと}
\begin{itemize}
  \item メモで見つけた「分け方の数が小さい」という構造は、そのまま計算量の削減に使える。
    制約が小さくて全探索で通るときも、メモの発想をコードに移せるかを一度考える。
  \item 分け目を並べるときは、両端（$k=0$ と $k=7$）を最初に書く。
  \item 区間 $[i,j)$ を全部試すときは、\verb|range| の終わりが「空の区間」「末尾まで」を含むか確かめ、端の入力（\verb|keyencex|）で試してから出す。
\end{itemize}

\end{document}
